%%=============================================================================
%% Voorwoord
%%=============================================================================

\chapter*{\IfLanguageName{dutch}{Woord vooraf}{Preface}}%
\label{ch:voorwoord}

%% TODO:
%% Het voorwoord is het enige deel van de bachelorproef waar je vanuit je
%% eigen standpunt (``ik-vorm'') mag schrijven. Je kan hier bv. motiveren
%% waarom jij het onderwerp wil bespreken.
%% Vergeet ook niet te bedanken wie je geholpen/gesteund/... heeft

% De keuze voor dit onderwerp was snel gemaakt toen ik van de vraag te weten kreeg. Het combineert analyseren van kaarten met de inzichten die data kunnen geven en duurzame energie, allemaal dingen die mij interesseren. De kans krijgen om bij te kunnen dragen aan een invloedrijk project van het formaat van de Einstein Telescoop maakte het des te beter. Om deze thesis tot stand te brengen hebben verschillende mensen hun tijd en kennis tot mijn beschikking gesteld en ik wens deze mensen dan ook te bedanken hiervoor.
% Allereerst wens ik mijn promotor de Hogent Veerle Depestele te bedanken voor haar hulp in het correct formuleren van mijn onderzoeksvraag en haar kritische blik hierop. Het heeft deze thesis beter en steker gemaakt. Ook verdiend ze een bedankje om begrip te tonen voor de drukke situatie waarin ik verkeer en niet te hard te zijn op deadlines.

% Een tweede bedankje moet uitgaan naar mijn co-promotor Prof. dr. Bart Vermang en zijn college dr. Robbe Breugelmans. Ze zijn zeer vrijgevig geweest met hun kostbare tijd en hun kennis. Ook wil ik hen bedanken om mij de kans te geven mee te werken aan deze unieke ervaring.

% Tot slot wil ik ook mijn vrouw Hannah Loots bedanken voor de ondersteuning in zowel het schrijven en nalezen van deze thesis alsook de extra ondersteuning in ons huishouden waardoor dit alles mogelijk is. 

% Ik wil ook graag mijn vader Willy Ulenaers bedanken om de tijd te nemen deze thesis taalkundig te controleren 

Kaarten hebben me altijd gefascineerd, niet alleen als afbeeldingen maar ook als verhalen die wachten om ontdekt te worden. Als jongvolwassene ontdekte ik de kracht van data en hoe het verborgen patronen blootlegt en inzichten biedt die je niet zomaar ziet. Ook is er mijn achtergrond in duurzame energie, een domein dat me blijft boeien, zeker als het gaat om innovatieve ontwikkelingen en hun toepassingen.
Dit project brengt al die interesses samen. Toen ik het onderwerp voor het eerst tegenkwam, was ik meteen geïntrigeerd. De kans om bij te dragen aan een groot onderzoeksproject in mijn eigen omgeving maakte het alleen maar aantrekkelijker.\medskip

Om deze thesis tot stand te brengen, hebben verschillende mensen hun tijd en kennis ter beschikking gesteld. Ik wil hen hiervoor bedanken.\medskip

Allereerst bedank ik mijn promotor, Veerle Depestele (HoGent), voor haar hulp bij het correct formuleren van mijn onderzoeksvraag en haar kritische blik. Haar feedback heeft deze thesis sterker gemaakt. Daarnaast waardeer ik haar begrip voor mijn drukke situatie en de flexibiliteit rondom deadlines.\medskip

Een groot dankwoord gaat uit naar mijn copromotor, Prof.\ dr.\ Bart Vermang (UHasselt), en zijn collega dr. Robbe Breugelmans. Jullie hebben vrijgevig jullie kostbare tijd en kennis gedeeld en mij de unieke kans gegeven om mee te werken aan dit project.\medskip

Tot slot wil ik mijn vrouw, Hannah, bedanken voor haar ondersteuning bij het schrijven en nalezen van deze thesis, alsook voor de extra inzet in ons huishouden, waardoor dit alles mogelijk was. Ook bedank ik mijn vader, Willy, voor het taalkundig controleren van de tekst.\medskip
